% TOP_PROBLEM: {"id": 1260, "title": "Is 10 a Solitary Number?", "category_id": 1, "category": {"id": 1, "name": "number_theory", "display_name": "Number Theory", "description": "Properties of integers, prime numbers, Diophantine equations.", "slug": "number-theory", "order_index": 1, "created_at": "2026-07-31T15:26:25.670Z"}, "status": "open"}
% ATTEMPT_STATUS: unresolved
% Research attempt by GPT_6_Astra_Ultra, 1 October 2026.
% Canonical UnsolvedMath ID; original statements and sources preserved.
% FIRST_POSTED: 2026-10-01
% Imported from: unsolvedmath_additional_500.tex; UnsolvedMath ID 1260
% !TeX program = lualatex
% Compile twice: lualatex 1260.tex
% Self-contained: no external bibliography, images, or \input files.
% Numeric section labels and visible numbers are original UnsolvedMath IDs.
\documentclass[11pt,a4paper]{article}
\usepackage[margin=24mm,headheight=14pt]{geometry}
\usepackage{fontspec}
\setmainfont{Latin Modern Roman}
\setsansfont{Latin Modern Sans}
\setmonofont{Latin Modern Mono}
\usepackage{amsmath,amssymb,mathtools}
\usepackage{unicode-math}
\setmathfont{Latin Modern Math}
\usepackage{microtype}
\usepackage{enumitem,xurl,needspace}
\usepackage[dvipsnames]{xcolor}
\usepackage{hyperref}
\hypersetup{unicode=true,colorlinks=true,linkcolor=MidnightBlue,urlcolor=MidnightBlue,
 pdftitle={UnsolvedMath Problem 1260},pdfauthor={GPT 6 Astra Ultra},bookmarksnumbered=true}
\usepackage{bookmark,tocloft,titlesec,fancyhdr}
\setlength{\cftsecnumwidth}{4.2em}
\cftsetpnumwidth{2.5em}
\cftsetrmarg{4em}
\titleformat{\section}{\Large\bfseries\raggedright}{\thesection}{0.7em}{}
\pagestyle{fancy}\fancyhf{}
\fancyhead[L]{\small\sffamily GPT 6 Astra Ultra research attempt}
\fancyhead[R]{\small Original catalogue IDs}
\fancyfoot[C]{\thepage}
\setlength{\parindent}{0pt}
\setlength{\parskip}{5pt plus 1pt}
\setlength{\emergencystretch}{3em}
\setcounter{tocdepth}{1}\setcounter{secnumdepth}{1}
\setlist[itemize]{leftmargin=3em,itemsep=4pt,topsep=4pt,parsep=0pt}
\allowdisplaybreaks
\newcommand{\N}{\mathbb{N}}
\newcommand{\Z}{\mathbb{Z}}
\newcommand{\Q}{\mathbb{Q}}
\newcommand{\R}{\mathbb{R}}
\newcommand{\C}{\mathbb{C}}
\newcommand{\F}{\mathbb{F}}
\newcommand{\A}{\mathbb{A}}
\newcommand{\PP}{\mathbb{P}}
\newcommand{\E}{\mathbb{E}}
\newcommand{\eps}{\varepsilon}
\newcommand{\dd}{\,\mathrm d}
\newcommand{\sref}[2]{\hyperlink{source:#1:#2}{[S#2]}}
\newif\ifshowcatalogue
\showcataloguetrue
\newcommand{\cataloguescope}[1]{\ifshowcatalogue
 \paragraph{Statement recorded in the supplied source.}
 #1\fi}
\begin{document}
% UnsolvedMath ID 1260; source catalogue code NT-053

\setcounter{section}{1259}

\section{Is Ten Solitary?}\label{1260}

{\small\sffamily UnsolvedMath ID 1260 · Number Theory}

\subsection{Definitions and mathematical statement}

\paragraph{Shared notation.}Unless a section says otherwise, $\N=\{1,2,\ldots\}$,
$\Z$ is the ring of integers, $\Q,\R,\C$ are the rational, real, and complex
fields, $[N]=\{1,\ldots,\lfloor N\rfloor\}$, and $\log$ is the natural
logarithm. For real functions of a parameter tending to infinity,
$f\ll g$ and $f=O(g)$ mean $|f|\leq Cg$ eventually for some fixed $C$;
$f\gg g$ reverses this comparison; $f=o(g)$ means $f/g\to0$;
$f\sim g$ means $f/g\to1$; and $f\asymp g$ means both order bounds.
A subscript on $O$ or $\ll$ specifies allowed constant dependence.
The expression $n^{o(1)}$ in an upper bound means growth smaller than
$n^\epsilon$ for each fixed $\epsilon>0$. Local definitions govern any
exception, finite-field convention, or use of the same symbol for another
object. An asymptotic claim is not automatically an effective algorithm.



The abundancy index is $I(n)=\sigma(n)/n$. A positive integer is solitary if no different positive integer has the same index. Since $\sigma(10)=18$, the question is whether $5\sigma(m)=9m$ has a positive integer solution other than $m=10$. \sref{1260}{1} \sref{1260}{2}

\paragraph{Statement recorded in the supplied source.}
Is 10 a solitary number (no other number shares its abundancy index)? \sref{1260}{1}

\paragraph{Residual task reported by the archive.}

The embedded review (2026-08-17) records: Find a friend of 10 or prove that sigma(n)/n=9/5 forces n=10. \sref{1260}{1}

\subsection{Short English statement}

Is ten the only positive integer whose sum of divisors is nine fifths of the integer?

\subsection{Research attempt}

\paragraph{Scope and source correction.}
The equation is $5\sigma(m)=9m$, and $m=10$ must be excluded as the
trivial solution. The URL in catalogue source [S2] identifies Henry
Robert Thackeray's paper, not the author imported in that citation.
The corrected primary reference is [E1]. Its stated ten-prime-factor
bound is not reproduced or needed below. The attempt proves a smaller
set of structural exclusions directly.

\paragraph{Every other candidate is an odd square divisible by five.}
The equation implies $5\mid m$. Abundancy $I(t)=\sigma(t)/t$ strictly
increases under a proper divisibility extension: prime by prime, each
additional exponent increases the finite sum $1+p^{-1}+\cdots+p^{-e}$,
and a new prime factor multiplies it by a factor greater than one.
If $m$ is even, then $10\mid m$, so $I(m)=I(10)$ forces $m=10$.
Thus every other candidate is odd.

For odd $m$, $\sigma(m)=9m/5$ is odd. The divisor-sum parity rule
therefore makes every prime exponent even, so $m$ is a square and
$5^2\mid m$. These conclusions cover all even candidates, with no
finite bound, and reduce the remaining search to odd squares.

\paragraph{The prime three can also be excluded.}
Suppose $3\mid m$. If $3^4\mid m$, then $3^4 5^2\mid m$ and
\[
 I(m)\ge\frac{121}{81}\frac{31}{25}
        =\frac{3751}{2025}>\frac95,
\]
contradiction. Since exponents are even, the only remaining option is
$3^2\Vert m$. Then $\sigma(3^2)=13$ divides $\sigma(m)$, and the
equation $5\sigma(m)=9m$ forces $13\mid m$. Hence
$3^2 5^2 13^2\mid m$, and
\[
 I(m)\ge\frac{13}{9}\frac{31}{25}\frac{183}{169}
        >\frac95.
\]
For the last strict comparison, cross multiplication reduces to
$13\cdot31\cdot183>405\cdot169$, namely $73749>68445$.
Thus $3\nmid m$. Every friend of ten would be an odd square coprime
to six, with five among its prime divisors.

\paragraph{A forced prime congruence and exponent pattern.}
Write $5^{2b}\Vert m$ with $b\ge1$. Taking five-adic valuations in
the equation gives $v_5(\sigma(m))=2b-1>0$. The factor
$\sigma(5^{2b})$ is one modulo five, so some other prime power
$q^{2e}\Vert m$ has $5\mid\sigma(q^{2e})$.
If $q\not\equiv1\pmod5$, the geometric-series identity implies
$q^{2e+1}\equiv1\pmod5$. The order of $q$ modulo five then divides
both the odd number $2e+1$ and four, so it is one, a contradiction.
Therefore $q\equiv1\pmod5$, and now
$\sigma(q^{2e})\equiv2e+1\pmod5$, yielding
\[
                   q\equiv1\pmod5,\qquad 2e\equiv4\pmod{10}.
\]
Thus at least one non-five prime occurs to exponent at least four.
The argument proves existence of such a factor, but does not identify
which factor or bound how large it can be.

\paragraph{Prime support must absorb divisor-sum factors.}
If $p^{2e}\Vert m$, every prime divisor of $\sigma(p^{2e})$ divides
$9m$, by the defining equation. Except for the prime three, such primes
must already belong to the support of $m$. At the same time $3\nmid m$
but $v_3(\sigma(m))=2$, so all contributions to the prime three from
the several divisor sums must total exactly two. Merely recording that
three occurs is insufficient; its valuation budget is fixed.

These conditions give a branching search: choose exact exponents,
factor their divisor sums, add forced primes, and prune branches with
abundancy exceeding $9/5$ or excessive three-adic valuation. The
unproved step is termination across arbitrarily many primes and
arbitrarily large even exponents. A contradiction in the all-exponents-two
branch does not address the forced higher exponent just derived.

\paragraph{Executed exact check.}
The script factors every $1\le m\le200000$, evaluates $\sigma(m)$,
and compares the integers $5\sigma(m)$ and $9m$. It finds exactly
$m=10$. No rounded abundancy values are compared. Code and result are
\texttt{checks\_second.py}, \texttt{results\_second.json} in
\path{_Local/Erdos_problems/UnsolvedMath/attempt_audit_2026_10_01/number_theory/}.
This bound is a modest local check, not the cited paper's structural
computation or a claim that every possible friend has been exhausted.

\paragraph{Remaining gap and outcome.}
All nontrivial candidates are reduced to odd squares divisible by five,
coprime to six, with a specified higher-exponent prime congruence and
exact divisor-sum valuation constraints. None of these conditions is
shown globally incompatible. \textbf{Outcome: UNRESOLVED.} No friend
of ten is constructed and no proof that ten is solitary is obtained.

\subsection{Sources}

{\small

\begin{itemize}

\item[\hypertarget{source:1260:1}{S1}] ULAM.AI, \emph{UnsolvedMath}, supplied \texttt{problems.json}, record 1260 (NT-053), statement and background. The embedded literature-review date is 2026-08-17; that is the archive’s review date, not a fresh certification. Dataset landing page: \url{https://huggingface.co/datasets/ulamai/UnsolvedMath}.

\item[\hypertarget{source:1260:2}{S2}] S. R. L. S. de Araujo, Each friend of 10 has at least 10 nonidentical prime factors, arXiv:2310.15900. \url{https://arxiv.org/abs/2310.15900}. Bibliographic information imported from the supplied record.


\item[E1] Henry Robert Thackeray, \emph{Each friend of 10 has at
least 10 nonidentical prime factors}, arXiv:2310.15900v2 (2024).
\url{https://arxiv.org/abs/2310.15900v2}. Author metadata and abstract
inspected 1 October 2026. This corrects the attribution attached to the
same URL in [S2]; the imported citation is preserved as provenance.

\end{itemize}

}

\label{end:1260}
\end{document}
